% !TEX TS-program = pdflatexmk

\documentclass[14pt]{beamer}
\usepackage{newtxtext,newtxmath}
\usepackage{microtype}
\usepackage[english]{babel}
\usepackage{booktabs}
\usepackage{graphicx}
\usepackage{tikz}

% Define UA colors
% https://brand.arizona.edu/guide/color
\definecolor{ua-red}{HTML}{AB0520}
\definecolor{ua-blue}{HTML}{0C234B}

\mode<presentation>{
\usetheme{Madrid}
\usecolortheme[named=ua-red]{structure}
\setbeamertemplate{navigation symbols}{}
\setbeamertemplate{footline}[frame number]
\setbeamertemplate{section in toc}[square]
\setbeamertemplate{subsection in toc}[square]
\setbeamertemplate{items}[square]
\setbeamercovered{transparent=0}
}

\author[Bethard]{Dr. Steven Bethard}
\institute[Arizona]{%
School of Information\\
University of Arizona}

\AtBeginSection[]
{
  \begin{frame}<beamer>{Outline}
    \tableofcontents[currentsection]
  \end{frame}
}

\title{Adapting natural language processing models across clinical domains}
\date[]{26 Feb 2021}


\begin{document}


\begin{frame}
  \titlepage
\end{frame}

\section{Domain adaptation is critical for clinical NLP}

\begin{frame}{Machine learning needs adaptation}
\small
\begin{tikzpicture}[data/.style={font=\scriptsize}]

\node[data, label={Source Domain}] (source-data) at (0, 0) {
\begin{tabular}{@{} l l @{}}
\toprule
Input & Output \\
\midrule
Dec 5, 2007 & 2007-12-05 \\
3 Jan 2018 & 2018-01-03 \\
04-25-1990 & 1990-04-25 \\
\ldots & \ldots \\
\bottomrule
\end{tabular}
};

\node[label={Learned Model}, inner sep=0pt] (model) at (5, 0) {\includegraphics{Hey_Machine_Learning_Logo.png}};

\draw[ultra thick, -latex]  (source-data)  -- node[align=center] {Machine \\ Learning} (model);

\visible<2->{
\node[data, label={Target Domain}, anchor=south west] (target-data) at (7, 0.5) {
\begin{tabular}{@{} p{9em} @{}}
\toprule
Input \\
\midrule
3 weeks postoperative \\
1 tablet q.i.d \\
75-year-old \\
\ldots \\
\bottomrule
\end{tabular}
};
\draw[ultra thick, -latex]  (target-data.south west)  -- (model.east);
}

\visible<3->{
\node[data, anchor=north west] (output) at (7, -0.5) {
\begin{tabular}{@{} p{9em} @{}}
\toprule
Output \\
\midrule
??? \\
??? \\
??? \\
\ldots \\
\bottomrule
\end{tabular}
};

\draw[ultra thick, -latex]  (model.east)  -- (output.north west);
}
\end{tikzpicture}
\end{frame}

\begin{frame}{Machine learning needs adaptation}
A sample of recent NLP models, trained on a source domain, and tested on a target domain:

\bigskip
\footnotesize
\begin{tabular}{ l l l l l }
\toprule
& \multicolumn{2}{c}{Source} & \multicolumn{2}{c}{Target} \\
\cmidrule(lr){2-3}\cmidrule(lr){4-5}
Task & Domain & Score & Domain & Score \\
\midrule
Sentences & Mayo Clinic & 94.9 A & BIDMC ICU & 31.8 A \\
Syntax trees & Wall Street Journal & 89 $F_1$ & Medline abstracts & 75 $F_1$ \\
Negations & Partners Healthcare & 93.6 $F_1$ & Mayo Clinic & 74.7 $F_1$ \\
Times & Mayo colon cancer & 80 $F_1$ & Mayo brain cancer & 59 $F_1$ \\
\ldots \\
\bottomrule
\end{tabular}

\centering
\smallskip
Full table in \cite{laparra:JAMIAOpen:2020}
\end{frame}

\begin{frame}{Domain adaptive machine learning}
\small
\begin{tikzpicture}[data/.style={font=\scriptsize}]

\node[data, label={Source Domain}] (source-data) at (0, 2) {
\begin{tabular}{@{} l l @{}}
\toprule
Input & Output \\
\midrule
Dec 5, 2007 & 2007-12-05 \\
3 Jan 2018 & 2018-01-03 \\
04-25-1990 & 1990-04-25 \\
\ldots & \ldots \\
\bottomrule
\end{tabular}
};

\node[data, label={Target Domain}] (target-data) at (0, -2) {
\begin{tabular}{@{} l l @{}}
\toprule
Input & \color{lightgray} Output\\
\midrule
1 week postop &  \color{lightgray} 2011-07-06 \\
2 tablets b.i.d. &  \color{lightgray} 2$\times$2$\times$P1D \\
40-year-old &  \color{lightgray} P40Y \\
\ldots \\
\bottomrule
\end{tabular}
};

\node[label={Learned Model}, inner sep=0pt] (model) at (5, 0) {\includegraphics{Hey_Machine_Learning_Logo.png}};

\node[align=center] at (2.8, 0) {Machine \\ Learning};

\draw[ultra thick, -latex]  (source-data)  -- (model);
\draw[ultra thick, -latex]  (target-data)  -- (model);

\visible<2->{\node[data, anchor=south west] (input) at (7, 0.5) {
\begin{tabular}{@{} p{9em} @{}}
\toprule
Input \\
\midrule
3 weeks postoperative \\
1 tablet q.i.d. \\
75-year-old \\
\ldots \\
\bottomrule
\end{tabular}
};

\node[data, anchor=north west] (output) at (7, -0.5) {
\begin{tabular}{@{} p{9em} @{}}
\toprule
Output \\
\midrule
2011-07-20 \\
1$\times$4$\times$P1D \\
P70Y \\
\ldots \\
\bottomrule
\end{tabular}
};

\draw[ultra thick, -latex]  (input.south west)  -- (model.east);
\draw[ultra thick, -latex]  (model.east)  -- (output.north west);
}
\end{tikzpicture}
\end{frame}

\begin{frame}{Domain adaptation is hard}
% new phrases
% new distribution of labels for old phrases
\end{frame}

\begin{frame}{Clinical domain adaptation is even harder}
% differences across institutions
% data sharing constraints
% etc.
\end{frame}

\section{Domain adaptation comes in many forms}

\begin{frame}{Taxonomy of domain adaptation}
% from JAMIA Open
\end{frame}

\begin{frame}{Frustratingly easy domain adaptation}
% diagram of method
% why it doesn't work for Clinical NLP
\end{frame}

\begin{frame}{Self-training}
% maybe the tri-training paper?
% diagram of method
% why it doesn't work for Clinical NLP
\end{frame}

\section{Domain adaptation for clinical information extraction}

\begin{frame}{Source-free domain adaptation}
% summarize shared tasks: negation and time
\end{frame}

\begin{frame}{Source-free self-training}
% Xin's results
\end{frame}

\begin{frame}{Online active learning}
% any results from Xin
\end{frame}

\section{Domain adaptation for clinical concept normalization}

\begin{frame}{Clinical concept normalization}
\end{frame}

\begin{frame}{Normalization as vector-space search}
\end{frame}

\begin{frame}{Domain adaptation without retraining}
\end{frame}

\section*{Summary and future work}

\begin{frame}{Summary}
\end{frame}

\begin{frame}{Future work}
\end{frame}

\appendix

\begin{frame}[allowframebreaks]
        \frametitle{References}
        \bibliographystyle{apalike}
        \bibliography{domain-adaptation.bib}
\end{frame}


\end{document}
